\subsection*{先分开三种状态}

服务在事务内部已经把余额改成 7000，另一连接仍读到 10000，客户端又只看见超时。这三个观察可以同时成立。若把它们写进一列统一状态，往往会产生含义含混的 \code{processing} 或 \code{unknown}：它们究竟说明数据库已经保存了什么，还是说明一个观察者暂时知道什么？检验需要为每类状态指定拥有者与观察位置。

先只看同一 \(S\) 的已提交事实。尚无该身份的终态时，状态记为\textbf{未形成}；完整成功提交后，记为\textbf{成功}，关联 P1、O1、D1 和原结果；余额条件不成立并保存完整拒绝后，记为\textbf{余额不足}。未形成可以变为成功或余额不足；在当前保留期内，两种终态都保持原样。退款如需发生，应作为另一项动作记录，保留创建购买的终态。

事务内部还有自己的执行位置：已插入身份、已修改余额、已插入订单、已插入流水、已填写终态、正在请求提交。\code{pending} 只表示事务尚在填写记录；普通外部查询看不到这一未提交行。任何正常出口都把完整终态提交，异常出口则撤销整笔事务。协议允许的发布步骤只有完整终态提交，内部构造过程由事务隔离。

客户端的状态则可以是待发送、结果待确认、已保存原结果。结果待确认涵盖了两条不同历史：一条请求在到达服务前丢失，数据库仍未形成 \(S\)；另一条已提交 P1，响应丢失。两条历史给客户端相同的超时观察，要求恢复时分别首次执行或读取 P1。\code{unknown} 记录在观察者的待确认状态或本次尝试记录中；恢复者继续查询数据库的未形成、成功或拒绝事实，再选择对应步骤。

\begin{figure}[htbp]
\centering
\begin{tikzpicture}
\node[booknote,anchor=west] at (0,3.7) {已提交数据库事实};
\node[booknode] (abs) at (2,2.8) {未形成};
\node[bookresult] (ok) at (7,3.4) {成功：P1、O1、D1、原结果};
\node[booknode] (no) at (7,2.1) {拒绝：P1、原因、原结果};
\draw[bookarrow] (abs) -- (ok);
\draw[bookarrow] (abs) -- (no);
\node[booknote] at (7,1.45) {两个终态分别保留，没有相互覆盖的迁移};
\node[booknode,text width=10.7cm] at (5.5,.3) {事务内部：插入身份 → 修改余额 → 写订单与流水 → 填写终态\\完整提交后发布；任一异常撤销内部组};
\node[booknode,text width=10.7cm] at (5.5,-1.3) {客户端观察：待发送 → 结果待确认 → 已保存结果\\超时所见的 unknown 属于此层，恢复继续携带原 S};
\end{tikzpicture}
\caption{数据库事实、事务内部位置与客户端观察分别变化}
\label{fig:ch06-state-layers}
\end{figure}


R1 超时后出现 R2，两者是两次 HTTP 尝试，仍可以属于同一项逻辑购买。逻辑操作按完整 \(S\) 归组，数据库终态说明这一组最后形成了什么。请求 ID 标记 R1 或 R2，购买编号标记最后成立的 P1；客户端首次发送前保存的 \(S\)，把后来的尝试重新接回同一项购买。

这种分开还影响崩溃记录。A 日志里出现候选 P1、扣后 7000，随后事务回滚，最终恢复者 B 用候选 P2 提交。已提交事实只有 P2，客户端最终保存的也应是 P2。编号出现过、语句执行过与业务已经成立，是三个分别需要取证的判断。

\subsection*{把允许的处理写成完整规则}

入口先核验账户权限、动作、许可与期限，再解析输入。商品 \code{paper-001}、数量 1、币种 CNY 与报价 Q7 按 \code{purchase-input-v1} 解释；该版本规定允许的等价写法怎样归一，例如把约定可省略的数量补为 1。规范化输入就是由这套规则得到的可比较字段。格式或字段错误只结束当前尝试，不占用购买身份。首次成交所需的 Q7 价格依据已在本地获得且可验证，运行过程不调用独立的外部扣款系统。

进入数据库事务以后，受理动作是对完整身份的受约束插入。四列均为非空，组合唯一约束限制相同四元组最多形成一条购买。服务检查实际返回的行数，取得身份才走余额更新；返回零行则在下一条独立语句读取原记录。这个分支是控制流的必要部分。

\begin{BookCode}
INSERT INTO purchases (
  purchase_id, account_id, action, session_id,
  idempotency_key, normalization_version,
  canonical_input, status
) VALUES (
  :candidate_purchase, 1, 'purchase.create.v1', 'E1',
  'K1', 'purchase-input-v1', :canonical_input, 'pending'
)
ON CONFLICT (account_id, action, session_id, idempotency_key)
DO NOTHING
RETURNING purchase_id;
\end{BookCode}

SQL 中冒号表示由服务绑定的参数，\code{candidate\_purchase} 是尚未成立的候选购买号。\code{DO NOTHING} 使目标身份冲突走正常的零行分支，\code{RETURNING} 只给出真正插入的编号。若冲突来自另一个未结束事务，数据库等待其结局：它提交，当前插入放弃候选；它回滚，当前尝试才可能插入成功。

取得身份后执行带条件的余额更新，条件与减法在数据库处理账户行时一起完成。成功返回当次扣后余额，随后写订单、流水和完整结果；返回零行，在账户存在、权限已确认的前提下保存余额不足。两条分支都在身份所在事务内终结。

\begin{BookCode}
开始事务，受约束地插入完整 S 与规范化输入
插入一行：验证本次首次成交依据
    条件 UPDATE：余额足够才减去 3000
    成功：写订单、流水、成功终态及原结果
    余额不足：写拒绝终态及原结果，不写订单和扣款
    提交完整事务，明确确认后才返回新结果
插入零行：下一条 SELECT 按完整 S 读原记录
    按原规范化版本比较本次输入
    同参：返回保存的终态与原结果
    异参：返回冲突，原记录保持不变
事务中任一步异常：回滚；结果未确认则保留 S 恢复
\end{BookCode}

\begin{figure}[htbp]
\centering
\begin{tikzpicture}
\node[booknode,text width=8cm] (entry) at (5.5,4.2) {核验资格与输入；开始事务；按完整 S 受约束插入};
\node[booknode,text width=4.4cm] (new) at (2.6,2.5) {实际插入一行\\条件更新余额};
\node[booknode,text width=4.4cm] (old) at (8.4,2.5) {插入零行\\下一条独立 SELECT};
\draw[bookarrow] (entry) -- (new);
\draw[bookarrow] (entry) -- (old);
\node[bookresult,text width=4.4cm] (write) at (2.6,.5) {成功：P、O、D、原结果\\或拒绝：P、原因、原结果\\全部在同一事务提交};
\node[booknode,text width=4.4cm] (read) at (8.4,.5) {按原版本比较输入\\同参读取原结果\\异参返回 409，保留旧记录};
\draw[bookarrow] (new) -- (write);
\draw[bookarrow] (old) -- (read);
\node[booknote,text width=10.5cm] at (5.5,-1) {右侧只读取与比较；只有实际插入身份的左侧进入首次扣款};
\end{tikzpicture}
\caption{身份插入决定首次处理与结果恢复的分支}
\label{fig:ch06-admission}
\end{figure}


冲突后的读取遵循第 4.3 节的两条语句规则，不在同一条 WITH 语句中查找原行；所有首次扣款都以实际插入身份成功为入口。

首次成功返回 201，首次已提交余额不足继续返回 409；同参恢复两种终态均返回 200。语法错误为 \code{400 invalid\_input}，字段语义错误为 \code{422 invalid\_input}，只说明本次未受理；异参冲突也为 409，查询未找到为 404，暂时未能确认结果为 503。调用方结合业务码与已保存的操作资料处理：遇到 404 或 503，继续保留 \(S\)，按原身份查询或恢复；另一次明确确认的购买才分配新身份。

原记录读到以后，按第五章的 \code{purchase-result-v1} 输出保存的成功或拒绝。成功结果中的 \code{balance\_after\_minor} 来自原结果或原 D1；当前余额另行查询账户并用独立字段展示。

\subsection*{在每个写入边界实际终止进程}

为了让退出确实落在指定位置，可以在专用购买工作进程中设置屏障。屏障使它完成某一步后暂停，父进程收到到达信号，再终止该工作进程。数据库服务器继续运行，恢复从另一连接开始。这样可以精确检查应用在身份、扣款或结果写完后突然消失时，数据库是否留下未完成事实。

沿写入路径设置八个位置：身份写入前、身份写入后、余额更新后、订单写入后、流水写入后、终态写入后、提交前，以及提交确认后但购买结果尚未交给调用方。各次都从 10000 分、空记录开始。前七处终止时尚未提交，预期新连接读到余额 10000，购买、订单与流水为空；同 \(S\) 恢复后只留各一条成功记录与余额 7000。最后一处 P1 已成立，预期恢复读回原结果，余额继续为 7000。这样先写出预期，才能把中断后的实际记录逐项与协议比较。

\begin{figure}[htbp]
\centering
\begin{tikzpicture}
\node[booknode,text width=10.7cm] (pre) at (5.5,3.1) {1.身份前　2.身份后　3.余额后　4.订单后\\5.流水后　6.终态后　7.提交前};
\node[booknote] at (8.4,1.75) {七处均在提交以前终止工作进程};
\node[booknode,text width=4.6cm] (rollback) at (2.7,.8) {数据库撤销内部组\\余额 10000；P、O、D 为空};
\node[bookresult,text width=4.6cm] (restore) at (8.4,.8) {同 S 首次恢复\\各一份 P、O、D；余额 7000};
\draw[bookarrow] (pre.south) -- (rollback.north);
\draw[bookarrow] (rollback) -- (restore);
\node[booknode,text width=4.6cm] (post) at (2.7,-1.2) {8.提交确认后、交付前退出\\完整成功；余额 7000};
\node[bookresult,text width=4.6cm] (replay) at (8.4,-1.2) {同 S 读取原结果\\P、O、D 与余额保持};
\draw[bookarrow] (post) -- (replay);
\end{tikzpicture}
\caption{八个终止位置按提交边界分为共同回滚与完整恢复}
\label{fig:ch06-kill-boundaries}
\end{figure}


表中把事务内部位置与恢复前事实并排，便于核对中断到底取消了哪些内容。各行都从独立的 10000 分、空记录开始，避免把多次实验累计到同一余额。

\begin{table}[htbp]
\centering\small
\begin{tabularx}{\linewidth}{@{}p{3.0cm}XXX@{}}
\toprule
终止位置 & 新连接恢复前读取 & 同 \(S\) 恢复路径 & 恢复后事实\\
\midrule
身份插入前或后 & 无终态，余额 10000 & 插入身份后首次处理 & 一份成功，余额 7000\\
余额或订单写入后 & 无订单扣款，余额 10000 & 原未提交修改已撤销 & 一份成功，余额 7000\\
流水或终态写入后 & 无终态，余额 10000 & 重新取得身份 & 一份成功，余额 7000\\
提交之前 & 无终态，余额 10000 & 重新取得身份 & 一份成功，余额 7000\\
提交确认后响应前 & 完整成功，余额 7000 & 读取原业务结果 & 原记录保持，余额 7000\\
\bottomrule
\end{tabularx}
\caption{应用在八个边界退出后的已提交事实与恢复}
\end{table}

核对时同时读取 \(S\) 的数量、终态、订单和流水归属、报价所需金额、当次余额、原结果编号，以及初始余额减去全部已提交扣款是否等于账户余额。只看余额 7000，会漏掉一条错指订单的流水；只看一条购买，也会漏掉结果中的错误编号。还可以在余额更新后暂停原进程，让三个同身份恢复者同时到达，再终止原进程。三个等待者应只有一个首次提交，另两个读回同一份结果，预期来自唯一身份与共同事务的规则。

这套安排观察应用退出引起的连接结束、未提交事务撤销与已提交结果恢复，数据库服务器在过程中继续运行。把各位置读到的记录与预期逐项比较，才能确认中断以后是否仍沿用同一协议。断电、磁盘失效和真实网络中的 COMMIT 确认丢失需要各自的故障安排；下节先推演提交确认未到达时的三个数据库分支。
\FloatBarrier
